% GENERATED FILE. Edit canonical lessons, then run npm run generate:books.
% canonical-source-hash: fnv1a64:e00c96800e94d497
% canonical-lessons: GU-C529-tabakko, GU-C529-kram, GU-C529-andaj, GU-C529-adhuru, GU-C529-antim

\chapter{One Stage at a Time}
\label{ch:gu-one-stage-at-a-time}
\hlchaptermodality{529}

\begin{chapteropening}
\textbf{By the end of this chapter:} \emph{I can say what happens if a stage goes well or is left half done.}

Complete the last lesson of chapter 529: I can say what happens if a stage goes well or is left half done.
\end{chapteropening}

\section[\texorpdfstring{\gu{તબક્કો} (tabakko)}{તબક્કો (tabakko)}]{\gu{તબક્કો} (tabakko) — a stage (of a process)}

\label{lesson:GU-C529-tabakko}



\begin{quote}
Before the new one: say the Gujarati for a diary, then the Gujarati for to organise.
\end{quote}

\subsection*{You'll want to know: \gu{તબક્કો}}
\textbf{\gu{તબક્કો}} — \emph{tabakko} — \textquotedblleft{}a stage (of a process)\textquotedblright{}.

\begin{quote}
\gu{જો} \gu{પહેલો} \gu{તબક્કો} \gu{સારો} \gu{જશે} \gu{તો} \gu{આપણે} \gu{બીજો} \gu{તબક્કો} \gu{શરૂ} \gu{કરીશું}.
\end{quote}

\begin{quote}
\emph{jo pahelo tabakko sāro jaśe to āpaṇe bījo tabakko śarū karīśũ.}
\end{quote}

\begin{quote}
If the first stage goes well, we will start the second stage.
\end{quote}

A masculine noun.

\begin{cousinweb}
\textbf{\gu{તબક્કો}} (\emph{tabakko}) is usually traced to Arabic \emph{ṭabaqa}, \textquotedblleft{}a layer, a stage\textquotedblright{}, through Persian.
\end{cousinweb}

\begin{grammarlens}[title={What you've built}]
One more word for saying what happens if.
\end{grammarlens}

\subsection*{Your turn}
\begin{itemize}
\raggedright
  \item \emph{Say it:} \emph{tabakko}
  \item \emph{Say it:} \emph{tabakko}, once more
  \item \emph{Say it:} \emph{jo pahelo tabakko sāro jaśe to āpaṇe bījo tabakko śarū karīśũ.} — If the first stage goes well, we will start the second stage.
  \item \emph{Recall:} say the Gujarati for a diary, then the Gujarati for to organise, then say \emph{tabakko} again
\end{itemize}

\begin{tcolorbox}[breakable,colback=teal!4,colframe=teal!35!black,title={Before you move on}]
What does \gu{તબક્કો} mean? (\textquotedblleft{}A stage (of a process)\textquotedblright{}.) Say it once more.
\end{tcolorbox}
\section[\texorpdfstring{\gu{ક્રમ} (kram)}{ક્રમ (kram)}]{\gu{ક્રમ} (kram) — a sequence}

\label{lesson:GU-C529-kram}



\begin{quote}
Before the new one: say the Gujarati for to organise, then the Gujarati for a stage (of a process).
\end{quote}

\subsection*{You'll want to know: \gu{ક્રમ}}
\textbf{\gu{ક્રમ}} — \emph{kram} — \textquotedblleft{}a sequence\textquotedblright{}.

\begin{quote}
\gu{જો} \gu{તમે} \gu{આ} \gu{ક્રમ} \gu{બદલશો} \gu{તો} \gu{બધું} \gu{બગડી} \gu{જશે}.
\end{quote}

\begin{quote}
\emph{jo tame ā kram badalśo to badhũ bagḍī jaśe.}
\end{quote}

\begin{quote}
If you change this order, everything will go wrong.
\end{quote}

A masculine noun.

\begin{cousinweb}
\textbf{\gu{ક્રમ}} (\emph{kram}) is a tatsama, taken from Sanskrit \emph{krama}, \textquotedblleft{}a step; going, a course\textquotedblright{} (Monier-Williams): one step after another.
\end{cousinweb}

\begin{grammarlens}[title={What you've built}]
One more word for saying what happens if.
\end{grammarlens}

\subsection*{Your turn}
\begin{itemize}
\raggedright
  \item \emph{Say it:} \emph{kram}
  \item \emph{Say it:} \emph{kram}, once more
  \item \emph{Say it:} \emph{jo tame ā kram badalśo to badhũ bagḍī jaśe.} — If you change this order, everything will go wrong.
  \item \emph{Recall:} say the Gujarati for to organise, then the Gujarati for a stage (of a process), then say \emph{kram} again
\end{itemize}

\begin{tcolorbox}[breakable,colback=teal!4,colframe=teal!35!black,title={Before you move on}]
What does \gu{ક્રમ} mean? (\textquotedblleft{}A sequence\textquotedblright{}.) Say it once more.
\end{tcolorbox}
\section[\texorpdfstring{\gu{અંદાજ} (andāj)}{અંદાજ (andāj)}]{\gu{અંદાજ} (andāj) — an estimate}

\label{lesson:GU-C529-andaj}



\begin{quote}
Before the new one: say the Gujarati for a stage (of a process), then the Gujarati for a sequence.
\end{quote}

\subsection*{You'll want to know: \gu{અંદાજ}}
\textbf{\gu{અંદાજ}} — \emph{andāj} — \textquotedblleft{}an estimate\textquotedblright{}.

\begin{quote}
\gu{જો} \gu{મારો} \gu{અંદાજ} \gu{સાચો} \gu{હશે} \gu{તો} \gu{આ} \gu{કામમાં} \gu{બે} \gu{દિવસ} \gu{લાગશે}.
\end{quote}

\begin{quote}
\emph{jo māro andāj sāco haśe to ā kāmmā̃ be divas lāgśe.}
\end{quote}

\begin{quote}
If my estimate is right, this job will take two days.
\end{quote}

A masculine noun.

\begin{cousinweb}
\textbf{\gu{અંદાજ}} (\emph{andāj}) is usually traced to Persian \emph{andāza}, \textquotedblleft{}a measure, an estimate\textquotedblright{}.
\end{cousinweb}

\begin{grammarlens}[title={What you've built}]
One more word for saying what happens if.
\end{grammarlens}

\subsection*{Your turn}
\begin{itemize}
\raggedright
  \item \emph{Say it:} \emph{andāj}
  \item \emph{Say it:} \emph{andāj}, once more
  \item \emph{Say it:} \emph{jo māro andāj sāco haśe to ā kāmmā̃ be divas lāgśe.} — If my estimate is right, this job will take two days.
  \item \emph{Recall:} say the Gujarati for a stage (of a process), then the Gujarati for a sequence, then say \emph{andāj} again
\end{itemize}

\begin{tcolorbox}[breakable,colback=teal!4,colframe=teal!35!black,title={Before you move on}]
What does \gu{અંદાજ} mean? (\textquotedblleft{}An estimate\textquotedblright{}.) Say it once more.
\end{tcolorbox}
\section[\texorpdfstring{\gu{અધૂરું} (adhūrũ)}{અધૂરું (adhūrũ)}]{\gu{અધૂરું} (adhūrũ) — half-finished}

\label{lesson:GU-C529-adhuru}



\begin{quote}
Before the new one: say the Gujarati for a sequence, then the Gujarati for an estimate.
\end{quote}

\subsection*{You'll want to know: \gu{અધૂરું}}
\textbf{\gu{અધૂરું}} — \emph{adhūrũ} — \textquotedblleft{}half-finished\textquotedblright{}.

\begin{quote}
\gu{જો} \gu{કામ} \gu{અધૂરું} \gu{રહેશે} \gu{તો} \gu{સાહેબ} \gu{નારાજ} \gu{થશે}.
\end{quote}

\begin{quote}
\emph{jo kām adhūrũ raheśe to sāheb nārāj thaśe.}
\end{quote}

\begin{quote}
If the work is left half-finished, the boss will be displeased.
\end{quote}

\begin{cousinweb}
\textbf{\gu{અધૂરું}} (\emph{adhūrũ}) comes from \emph{*ardhapūraka}, \textquotedblleft{}half-filled\textquotedblright{}, a form CDIAL reconstructs, and CDIAL lists the Gujarati word as \textquotedblleft{}half-done\textquotedblright{} (CDIAL 664).
\end{cousinweb}

\begin{grammarlens}[title={What you've built}]
One more word for saying what happens if.
\end{grammarlens}

\subsection*{Your turn}
\begin{itemize}
\raggedright
  \item \emph{Say it:} \emph{adhūrũ}
  \item \emph{Say it:} \emph{adhūrũ}, once more
  \item \emph{Say it:} \emph{jo kām adhūrũ raheśe to sāheb nārāj thaśe.} — If the work is left half-finished, the boss will be displeased.
  \item \emph{Recall:} say the Gujarati for a sequence, then the Gujarati for an estimate, then say \emph{adhūrũ} again
\end{itemize}

\begin{tcolorbox}[breakable,colback=teal!4,colframe=teal!35!black,title={Before you move on}]
What does \gu{અધૂરું} mean? (\textquotedblleft{}Half-finished\textquotedblright{}.) Say it once more.
\end{tcolorbox}
\section[\texorpdfstring{\gu{અંતિમ} (antim)}{અંતિમ (antim)}]{\gu{અંતિમ} (antim) — final}

\label{lesson:GU-C529-antim}



\begin{quote}
Before the new one: say the Gujarati for an estimate, then the Gujarati for half-finished.
\end{quote}

\subsection*{You'll want to know: \gu{અંતિમ}}
\textbf{\gu{અંતિમ}} — \emph{antim} — \textquotedblleft{}final\textquotedblright{}.

\begin{quote}
\gu{જો} \gu{અંતિમ} \gu{દિવસે} \gu{વરસાદ} \gu{પડશે} \gu{તો} \gu{આપણે} \gu{મેળો} \gu{બંધ} \gu{રાખીશું}.
\end{quote}

\begin{quote}
\emph{jo antim divse varsād paḍśe to āpaṇe meḷo bandh rākhīśũ.}
\end{quote}

\begin{quote}
If it rains on the final day, we will keep the fair closed.
\end{quote}

\begin{cousinweb}
\textbf{\gu{અંતિમ}} (\emph{antim}) is a tatsama, taken from Sanskrit \emph{antima}, \textquotedblleft{}last\textquotedblright{} (Monier-Williams), built on \emph{anta}, \textquotedblleft{}the end\textquotedblright{}, the word behind \textbf{\gu{અંત}}.
\end{cousinweb}

\begin{grammarlens}[title={What you've built}]
One more word for saying what happens if.
\end{grammarlens}

\subsection*{Your turn}
\begin{itemize}
\raggedright
  \item \emph{Say it:} \emph{antim}
  \item \emph{Say it:} \emph{antim}, once more
  \item \emph{Say it:} \emph{jo antim divse varsād paḍśe to āpaṇe meḷo bandh rākhīśũ.} — If it rains on the final day, we will keep the fair closed.
  \item \emph{Recall:} say the Gujarati for an estimate, then the Gujarati for half-finished, then say \emph{antim} again
\end{itemize}

\begin{tcolorbox}[breakable,colback=teal!4,colframe=teal!35!black,title={Before you move on}]
What does \gu{અંતિમ} mean? (\textquotedblleft{}Final\textquotedblright{}.) Say it once more.
\end{tcolorbox}
